\documentclass[10pt, letterpaper]{article}

% Packages:
\usepackage[
    ignoreheadfoot, % set margins without considering header and footer
    top=2 cm, % seperation between body and page edge from the top
    bottom=2 cm, % seperation between body and page edge from the bottom
    left=2 cm, % seperation between body and page edge from the left
    right=2 cm, % seperation between body and page edge from the right
    footskip=1.0 cm, % seperation between body and footer
    % showframe % for debugging 
]{geometry} % for adjusting page geometry
\usepackage{titlesec} % for customizing section titles
\usepackage{tabularx} % for making tables with fixed width columns
\usepackage{array} % tabularx requires this
\usepackage[dvipsnames]{xcolor} % for coloring text
\definecolor{primaryColor}{RGB}{0, 79, 144} % define primary color
\usepackage{enumitem} % for customizing lists
\usepackage{fontawesome5} % for using icons
\usepackage{amsmath} % for math
\usepackage[
    pdftitle={Omar ABEDELKADER's CV},
    pdfauthor={Omar ABEDELKADER},
    pdfcreator={LaTeX with RenderCV},
    colorlinks=true,
    urlcolor=primaryColor
]{hyperref} % for links, metadata and bookmarks
\usepackage[pscoord]{eso-pic} % for floating text on the page
\usepackage{calc} % for calculating lengths
\usepackage{bookmark} % for bookmarks
\usepackage{lastpage} % for getting the total number of pages
\usepackage{changepage} % for one column entries (adjustwidth environment)
\usepackage{paracol} % for two and three column entries
\usepackage{ifthen} % for conditional statements
\usepackage{needspace} % for avoiding page brake right after the section title
\usepackage{iftex} % check if engine is pdflatex, xetex or luatex
\usepackage{graphicx}

% Ensure that generate pdf is machine readable/ATS parsable:
\ifPDFTeX
    \input{glyphtounicode}
    \pdfgentounicode=1
    % \usepackage[T1]{fontenc} % this breaks sb2nov
    \usepackage[utf8]{inputenc}
    \usepackage{lmodern}
\fi



% Some settings:
\AtBeginEnvironment{adjustwidth}{\partopsep0pt} % remove space before adjustwidth environment
\pagestyle{empty} % no header or footer
\setcounter{secnumdepth}{0} % no section numbering
\setlength{\parindent}{0pt} % no indentation
\setlength{\topskip}{0pt} % no top skip
\setlength{\columnsep}{0cm} % set column seperation
\makeatletter
\let\ps@customFooterStyle\ps@plain % Copy the plain style to customFooterStyle
\patchcmd{\ps@customFooterStyle}{\thepage}{
    \color{gray}\textit{\small Omar ABEDELKADER - Page \thepage{} of \pageref*{LastPage} }
}{}{} % replace number by desired string
\makeatother
\pagestyle{customFooterStyle}

\titleformat{\section}{\needspace{4\baselineskip}\bfseries\large}{}{0pt}{}[\vspace{1pt}\titlerule]

\titlespacing{\section}{
    % left space:
    -1pt
}{
    % top space:
    0.3 cm
}{
    % bottom space:
    0.2 cm
} % section title spacing

\renewcommand\labelitemi{$\circ$} % custom bullet points
\newenvironment{highlights}{
    \begin{itemize}[
        topsep=0.10 cm,
        parsep=0.10 cm,
        partopsep=0pt,
        itemsep=0pt,
        leftmargin=0.4 cm + 10pt
    ]
}{
    \end{itemize}
} % new environment for highlights

\newenvironment{highlightsforbulletentries}{
    \begin{itemize}[
        topsep=0.10 cm,
        parsep=0.10 cm,
        partopsep=0pt,
        itemsep=0pt,
        leftmargin=10pt
    ]
}{
    \end{itemize}
} % new environment for highlights for bullet entries


\newenvironment{onecolentry}{
    \begin{adjustwidth}{
        0.2 cm + 0.00001 cm
    }{
        0.2 cm + 0.00001 cm
    }
}{
    \end{adjustwidth}
} % new environment for one column entries

\newenvironment{twocolentry}[2][]{
    \onecolentry
    \def\secondColumn{#2}
    \setcolumnwidth{\fill, 4.5 cm}
    \begin{paracol}{2}
}{
    \switchcolumn \raggedleft \secondColumn
    \end{paracol}
    \endonecolentry
} % new environment for two column entries

\newenvironment{header}{
    \setlength{\topsep}{0pt}\par\kern\topsep\centering\linespread{1.5}
}{
    \par\kern\topsep
} % new environment for the header

\newcommand{\placelastupdatedtext}{% \placetextbox{<horizontal pos>}{<vertical pos>}{<stuff>}
  \AddToShipoutPictureFG*{% Add <stuff> to current page foreground
    \put(
        \LenToUnit{\paperwidth-2 cm-0.2 cm+0.05cm},
        \LenToUnit{\paperheight-1.0 cm}
    ){\vtop{{\null}\makebox[0pt][c]{
        \small\color{gray}\textit{Last updated in December 2024}\hspace{\widthof{Last updated in December 2024}}
    }}}%
  }%
}%

% save the original href command in a new command:
\let\hrefWithoutArrow\href

% new command for external links:
\renewcommand{\href}[2]{\hrefWithoutArrow{#1}{\ifthenelse{\equal{#2}{}}{ }{#2 }\raisebox{.15ex}{\footnotesize \faExternalLink*}}}


\begin{document}
    \newcommand{\AND}{\unskip
        \cleaders\copy\ANDbox\hskip\wd\ANDbox
        \ignorespaces
    }
    \newsavebox\ANDbox
    \sbox\ANDbox{}

    \placelastupdatedtext
    \begin{header}
        \textbf{\fontsize{24 pt}{24 pt}\selectfont Omar ABEDELKADER}

        \vspace{0.3 cm}

        \normalsize
        \mbox{{\color{black}\footnotesize\faMapMarker*}\hspace*{0.13cm}Lille, France}%
        \kern 0.25 cm%
        \AND%
        \kern 0.25 cm%
        \mbox{\hrefWithoutArrow{mailto:omar.abedelkader@inria.fr}{\color{black}{\footnotesize\faEnvelope[regular]}\hspace*{0.13cm}omar.abedelkader@inria.fr}}%
        \kern 0.25 cm%
        \AND%
        \kern 0.25 cm%
        \mbox{\hrefWithoutArrow{https://omarabedelkader.github.io/}{\color{black}{\footnotesize\faLink}\hspace*{0.13cm}omarabedelkader.fr}}%
    \end{header}

    \vspace{0.3 cm - 0.3 cm}

    \section{About Me}  
        \begin{onecolentry}
            I am an AI researcher with extensive experience in Artificial Intelligence and Software Engineering. \\
            Currently pursuing a Ph.D. in Computer Science at the University of Lille. \end{onecolentry}
        \vspace{0.2 cm}

    \section{Education}     
        \begin{twocolentry}{
            \textit{Oct. 2024 – Oct. 2027}}
            \textbf{University of Lille}
            \textit{Villeneuve-d'Ascq, Nord}

            \textit{Ph.D in Computer Science}
            
        \end{twocolentry}
        \vspace{0.10 cm}

        \begin{twocolentry}{
            \textit{Sept. 2022 – Sept. 2024}}
            \textbf{University of Lorraine}
            \textit{Nancy, M\&M}

            \textit{M.Sc. in Natural Language Processing}
        \end{twocolentry}
        \vspace{0.10 cm}

        \begin{twocolentry}{
            \textit{Sept. 2019 – July 2022}}
            \textbf{Lebanese University}
            \textit{Beirut, Lebanon}

            \textit{B.Sc. in Data Science}
        \end{twocolentry}

        \vspace{0.10 cm}

%%%%%%%%%%%%%%%%%%%%%%% EXPERIENCE
    \section{Experience}


        \begin{twocolentry}{
            \textit{Villeneuve d'Asq, Nord}    
                
            \textit{Oct. 2024 – Oct. 2027}}
                \textbf{AI Researcher}
                
                \textit{INRIA}
        \end{twocolentry} 

        \vspace{0.10 cm}
        \begin{onecolentry}
            \begin{highlights}
                \item Improving code completion and generation using LLMs, specifically targeting the Pharo programming language, which has limited training data. 
                \item Developing techniques for code completion, type inference, and deployment in Pharo’s IDE, with a focus on runtime performance.
            \end{highlights}
        \end{onecolentry}
        \vspace{0.2 cm}

        \begin{twocolentry}{
            \textit{Verneuil-en-Halatte, Oise}    
                
            \textit{Sept. 2023 – Sept. 2024}}
                \textbf{AI Engineer}
                
                \textit{INERIS}
        \end{twocolentry} 
        \vspace{0.10 cm}
        \begin{onecolentry}
            \begin{highlights}
                \item Developing “INERIS-IA,” a tool to classify textual documents based on INERIS’s strategic goals using ML and NLP techniques
                \item Create boolean queries for document retrieval, and improve corpus quality through document similarity and keyword extraction.
            \end{highlights}
        \end{onecolentry}
        \vspace{0.2 cm}


        \begin{twocolentry}{
            \textit{Villetaneuse, SSD}    
                
            \textit{June 2023 – Aug. 2023}}
                \textbf{Intern - AI Researcher}
                
                \textit{LIPN}
        \end{twocolentry}    
        \vspace{0.10 cm}
        \begin{onecolentry}
            \begin{highlights}
                \item Comparison of various sampling techniques for probabilistic planning, particularly in generating literary narratives.
                \item Evaluate different methods, including the Score Function Estimator (SFE) and more advanced techniques like Gumbel-Softmax, to assess their effectiveness in creating coherent and creative stories.
            \end{highlights}
        \end{onecolentry}

%%%%%%%%%%%%%%%%%%%%%%% PROJECTS
        \section{Projets}


        \begin{twocolentry}{
            \textit{Nancy, M\&M}    
                
            \textit{Sept. 2023 – Jan. 2024}}
                \textbf{Family Dynamics Analysis}
                
                \textit{IDMC}
        \end{twocolentry} 

        \vspace{0.10 cm}
        \begin{onecolentry}
            \begin{highlights}
                \item  This project analyzes family dynamics in France using Formal Concept Analysis (FCA). By leveraging data from the French national census, the study identifies patterns in family structures, single-parent households, childbearing decisions, and naturalization through marriage across five regional zones. Using advanced data mining techniques, the project explores relationships between demographic variables such as marital status, household composition, and nationality, offering insights into the cultural and social trends shaping French households.
            \end{highlights}
            \begin{itemize}
                \item \textbf{GitHub Repository}: \href{https://github.com/omarabedelkader/FCA}{GitHub}
                \item \textbf{Paper}: \href{https://github.com/omarabedelkader/FCA/blob/main/paper/paper.pdf}{Paper}
              \end{itemize}              
        \end{onecolentry}
        \vspace{0.2 cm}


        \begin{twocolentry}{
            \textit{Nancy, M\&M}    
                
            \textit{Sept. 2023 – Jan. 2024}}
                \textbf{Real-Fake Face detection}
                
                \textit{IDMC}
        \end{twocolentry} 

        \vspace{0.10 cm}
        \begin{onecolentry}
            \begin{highlights}
                \item This project demonstrates the process of building and training a neural network for classification tasks using machine learning frameworks. It walks through data preprocessing, model architecture definition, training, evaluation, and visualization of results, providing insights into model performance and prediction accuracy.
            \end{highlights}
            \begin{itemize}
                \item \textbf{GitHub Repository}: \href{https://github.com/omarabedelkader/real-fake-detection}{GitHub}
              \end{itemize}              
        \end{onecolentry}
        \vspace{0.2 cm}

        \begin{twocolentry}{
            \textit{Nancy,  M\&M}    
                
            \textit{Sept. 2022 – June. 2023}}
                \textbf{EESMACF}
                
                \textit{IDMC \& LORIA}
        \end{twocolentry} 
        \vspace{0.10 cm}
        \begin{onecolentry}
            \begin{highlights}
                \item This project focuses on evaluating the effectiveness of two NLP models—SBERT (Sentence-BERT) and MiniLM—in classifying analogies using the FrameNet dataset. FrameNet is a rich semantic database that captures the meanings of lexical units within specific contexts called frames. The objective is to identify valid and invalid analogies in FrameNet, leveraging the capabilities of SBERT and a fine-tuned MiniLM model. SBERT is used to create dense vector embeddings of sentences that retain semantic similarity, while MiniLM is trained and fine-tuned to classify analogies effectively. The project found that MiniLM achieved an impressive 99\% accuracy in distinguishing between valid and invalid analogies, outperforming SBERT, which had around 55\% accuracy.
            \end{highlights}
            \begin{itemize}
                \item \textbf{GitHub Repository}: \href{https://github.com/omarabedelkader/EESMACF}{GitHub}
                \item \textbf{Paper}: \href{https://github.com/omarabedelkader/EESMACF/blob/main/paper/paper.pdf}{Paper}
              \end{itemize}
              
        \end{onecolentry}
        \vspace{0.2 cm}


        \begin{twocolentry}{
            \textit{IDMC, M\&M}    
                
            \textit{Sept. 2022 – Jan. 2023}}
                \textbf{DeGatto}
                
                \textit{IDMC}
        \end{twocolentry}    
        \vspace{0.10 cm}
        \begin{onecolentry}
            \begin{highlights}
              \item The "DeGatto" project is a sentiment analysis framework designed for e-commerce, specifically focusing on women’s apparel reviews. Using a dataset sourced from Kaggle, which includes over 23,000 sentences and aspect-level annotations for material, size, design, and comfort, the project aims to support e-commerce businesses and customers by analyzing feedback at both the sentence and aspect levels. Various NLP, DL, and ML models were tested, including LSTM (BiLSTM), SVM, Logistic Regression, and Multinomial Naive Bayes, with BiLSTM achieving the best results in sentence-level analysis and LinearSVC performing well at aspect-level analysis. A visualization tool was developed using ReactJS and NodeJS, enabling users to view results as bar or pie charts. The project concluded that BiLSTM is the most suitable model for sentence-level sentiment analysis, while LinearSVC excels in aspect-level analysis, providing a robust framework for sentiment classification in e-commerce contexts.      
            \end{highlights}

            \begin{itemize}
                \item \textbf{GitHub Repository}: \href{https://github.com/omarabedelkader/DeGatto/tree/main}{GitHub}
                \item \textbf{Paper}: \href{https://github.com/omarabedelkader/DeGatto/blob/main/paper/DeGatto.pdf}{Paper}
              \end{itemize}                  
        \end{onecolentry}   

%%%%%%%%%%%%%%%%%%%%%%% SKILLS

    \section{Technologies \& Skills}     
        \begin{onecolentry}
            \textbf{Programming Languages:} Python, Pharo, Java, R, JavaScript, PHP, MySQL
        \end{onecolentry}
        \vspace{0.2 cm}

        \begin{onecolentry}
            \textbf{Machine Learning \& Deep Learning:} TensorFlow, Keras, Scikit-learn
        \end{onecolentry}
        \vspace{0.2 cm}

        \begin{onecolentry}
            \textbf{Natural Language Processing (NLP):} SpaCy, NLTK
        \end{onecolentry}
        \vspace{0.2 cm}

        \begin{onecolentry}
            \textbf{Big Data \& Distributed Systems:} Apache Spark, Hadoop
        \end{onecolentry}
        \vspace{0.2 cm}
        \begin{onecolentry}
            \textbf{Data Analysis \& Visualization:} Pandas, NumPy, Matplotlib, Seaborn
        \end{onecolentry}
        \vspace{0.2 cm}

        \begin{onecolentry}
            \textbf{Frameworks \& Tools:} Flaks, Django, MongoDB, Laravel, ReactJS, AWS, PyCharm, 
        \end{onecolentry}
        \vspace{0.2 cm}

        \begin{onecolentry}
            \textbf{Operating Systems:} MacOS, UNIX, Windows
        \end{onecolentry}
        \vspace{0.2 cm}

        \begin{onecolentry}
            \textbf{Scientific Computing:} SciPy, NumPy
        \end{onecolentry}
        \vspace{0.2 cm}

        \begin{onecolentry}
            \textbf{Research Areas:} Artificial Intelligence, Software Engineer
        \end{onecolentry}

%%%%%%%%%%%%%%%%%%%%%%% VOLUNTEERING
    \section{Volunteering}

    \begin{onecolentry}
        \textbf{Ministère de la Transition Écologique:} Contributor in Club IA et Transition écologique
    \end{onecolentry}
    \vspace{0.2 cm}

    \begin{onecolentry}
        \textbf{Conseil National du Numérique (CNNum):} Contributor in Café IA
    \end{onecolentry}
    \vspace{0.2 cm}

    \begin{onecolentry}
        \textbf{Croix-Rouge Française:} Volunteer First Aider
    \end{onecolentry}

    %%%%%%%%%%%%%%%%%%%%%%% HOBBIES
    \section{Hobbies}
        \begin{onecolentry}
            \textbf{Water Activities:} Swimming · Diving
        \end{onecolentry}
        \vspace{0.2 cm}

        \begin{onecolentry}
            \textbf{Outdoor and Adventure Sports:} Climbing · Skiing · Skydiving
        \end{onecolentry}
        \vspace{0.2 cm}
        
        \begin{onecolentry}
            \textbf{Creative and Artistic Activities:} Photography
        \end{onecolentry}
        \vspace{0.2 cm}

        \begin{onecolentry}
            \textbf{Precision Sports:} Boxing · Bowling · Golf
        \end{onecolentry}
        \vspace{0.2 cm}


    %%%%%%%%%%%%%%%%%%%%%%% LANGUAGE
    \section{Language}
        \begin{onecolentry}
            \textbf{French:} Native
        \end{onecolentry}
        \vspace{0.2 cm}

        \begin{onecolentry}
            \textbf{English:} B2
        \end{onecolentry}
        \vspace{0.2 cm}


%%%%%%%%%%%%%%%%%%%%%%% REFERENCES        
    \section{References}
    Available on Request


    %\vspace{ cm}
    \begin{figure}[th]
        \centering
        \includegraphics[width=2cm, height=2cm]{qr_code_cv.png}
    \end{figure}
    \begin{center}
            \textbf{Scan Me !}
    \end{center}

\end{document}